\section{Solving Systems of Linear Equations}

  \begin{definition}[Linear System of Equations]
    Fix a field $\mathbb{F}$. A \textbf{linear equation} with variables $x_1, x_2, \cdots, x_n$ is in the form 
    \begin{equation}
      a_1 x_1 + a_2 x_2 + a_3 x_3 + \cdots + a_n x_n = b
    \end{equation}
    where the \textbf{coefficients} $a_i$ and the \textbf{free term} $b$ belong to $\mathbb{F}$. If $b = 0$, then $(3)$ is called a \textbf{homogeneous equation} and if $b \neq 0$, then it is called a \textbf{inhomogeneous equation}. 
  \end{definition}

  A system of $m$ linear equations with $n$ variables has the following general form 
  \begin{align*}
    &a_{1 1} x_1 + a_{1 2} x_2 + \cdots + a_{1 n} x_n = b_1 \\
    &a_{2 1} x_1 + a_{2 2} x_2 + \cdots + a_{2 n} x_n = b_2 \\
    & \cdots = \cdots\\
    &a_{m 1} x_1 + a_{m 2} x_2 + \cdots + a_{m n} x_n = b_m
  \end{align*}
  By matrix multiplication, this system is equal to the matrix equation $A x = b$.
  \begin{equation}
    \begin{pmatrix}
     a_{1 1} & a_{1 2} & \ldots & a_{1 n} \\
     a_{2 1} & a_{2 2} & \ldots & a_{2 n} \\
     \vdots & \vdots & \ddots & \vdots \\
     a_{m 1} & a_{m 2} & \ldots & a_{m n} 
    \end{pmatrix} \begin{pmatrix}
     x_1 \\ x_2 \\ x_3 \\ \vdots \\ x_n
    \end{pmatrix} = \begin{pmatrix}
    b_1 \\ b_2 \\ \vdots \\ b_m
    \end{pmatrix}
  \end{equation}
  That is, given a linear transformation $A: \mathbb{F}^n \longrightarrow \mathbb{F}^m$ and a vector $b \in \mathbb{F}^m$, we must find the preimage of $b$ under $A$. Clearly, $x$ is a solution of this matrix equation if and only if it is a solution of the system of equations. 

  We can interpret this matrix equation in two ways. First, we introduce the \textit{hyperplane interpertation}. The solution to each linear equation of $n$ variables represents an affine hyperplane in $\mathbb{F}^n$. Therefore, the solutions to the system of $m$ linear equations is simply the intersection of the $m$ affine hyperplanes of dimension $n-1$ within $\mathbb{R}^n$. That is, $x$ is a solution of $A x = b$ if and only if  
  \begin{equation}
    x \in \bigcap_{i = 1}^m \Big\{ (x_1, x_2, \cdots, x_n) \; | \; \sum_{j = 1}^n a_{i j} x_j = b_i \Big\}
  \end{equation}
  The \textit{column space interpretation} presents $A x = b$ in this equivalent form. 
  \begin{equation}
    x_1 \begin{pmatrix}
    a_{1 1} \\ a_{2 1} \\ \vdots \\ a_{m 1}
    \end{pmatrix} + x_2 \begin{pmatrix}
    a_{1 2} \\ a_{2 2} \\ \vdots \\ a_{m 2}
    \end{pmatrix} + \ldots + x_n \begin{pmatrix}
    a_{1 n} \\ a_{2 n} \\ \vdots \\ a_{m n}
    \end{pmatrix} = \begin{pmatrix}
    b_1 \\ b_2 \\ \vdots \\ b_m
    \end{pmatrix}
  \end{equation}
  That is, the solutions $x_1, x_2, \cdots, x_n$ are precisely the coefficients of the linear combination of the column vectors of $A$ that add up to vector $b$. Equivalently, it is the realization of vector $b$ with respect to the coordinate system of the column vectors of $A$. Note that the column space need not be a basis of $\mathbb{F}^m$. It does not need to be linearly independent nor does it need to span $\mathbb{F}^n$. 

  \begin{definition}[Coefficient Matrix]
    The matrix $A$ under the system is called the \textbf{coefficient matrix} and the matrix 
    \begin{equation}
      \Tilde{A} \equiv \begin{pmatrix}
      | & | & \cdots & | & | \\
      a_1 & a_2 & \cdots & a_n & b \\
      | & | & \cdots & | & | 
      \end{pmatrix} \equiv \begin{pmatrix}
      a_{1 1} & a_{1 2} & \ldots& a_{1 n} & b_1\\
       a_{2 1} & a_{2 2} & \ldots & a_{2 n} & b_2\\
      \vdots & \vdots & \ddots & \vdots & \vdots\\
       a_{m 1} & a_{m 2} & \ldots & a_{m n} & b_m
      \end{pmatrix}
    \end{equation}
    is called the \textbf{extended matrix}. 
  \end{definition}

  Clearly, these two definitions are equivalent since every elementary transformation of a system has a corresponding row transformation in its extended matrix. Under the hyperplane interpretation, it is a bit harder to visualize, but it is sufficient to say that each elementary operation either "stretches/compresses" (iii) a hyperplane or it "rotates" (i) the hyperplane around the axis where the solution exists. Either way, the intersection between the hyperplane and the set of solutions do not change. 

  \begin{definition}
    A system of linear equations is said to be \textbf{solved} if its extended matrix is in reduced row echelon form. 
  \end{definition}

  \begin{theorem}
    Elementary operations on either a system of linear equations or its extended matrix does not change its solutions. 
  \end{theorem}
  \begin{proof}
    It is easy to see this is true when performing the computations with the three transformations. We can prove this more abstractly (tbd): Given the system $A x = b$ with $x \in \mathbb{F}^n, b \in \mathbb{F}^m$. We see that $A \in $ Mat$(m \times n, \mathbb{F}) \implies \Tilde{A} \in$ Mat$(m \times (n+1), \mathbb{F})$. Each elementary row transformation on $\Tilde{A}$, denote it $E$, is a bijective mapping. Let us define the mapping 
    \begin{equation}
      \text{sol}: \text{Mat}\big( m \times (n+1), \mathbb{F} \big) \longrightarrow 2^{\mathbb{F}^n}, \; \text{sol} \begin{pmatrix}
      A & b
      \end{pmatrix} \equiv \{ x\in \mathbb{F}^n \; | \; A x = b\} 
    \end{equation}
    where $2^{\mathbb{F}^n}$ is the set of all subsets of $\mathbb{F}^n$. By matrix multiplication, we see that 
    \begin{equation}
       E \begin{pmatrix}
      A & b
      \end{pmatrix} = \begin{pmatrix}
      E A & E b
      \end{pmatrix}
    \end{equation}
    Since $E$ is bijective, it is invertible. So, 
    \begin{align} 
      \text{sol} \big( E \begin{pmatrix}
      A & b \end{pmatrix} \big) & = \text{sol} \begin{pmatrix}
      E A & E b \end{pmatrix} \\ 
      & = \{ x \mid E A x = E b\} \\ 
      & = \{ x \mid A x = b\} \\ 
      & = \text{sol} \begin{pmatrix} A & b \end{pmatrix}
    \end{align}
  \end{proof}

  Note the importance of this theorem. This result is the foundation behind the applications of Jordan Elimination.

  \begin{definition}
    A linear system can have either have no possible solutions (\textbf{overdetermined}), one unique solution, or multiple solutions (\textbf{underdetermined}) (infinite solutions if char$\, \mathbb{F}$ = 0). We can say with probability 1 that given a random $m \times n$ matrix $A$ with random $m$-dimensional vector $b$, the system $A x = b$ has
    \begin{enumerate}
      \item 0 solutions if $m > n$, since there are more equations than variables
      \item 1 solution if $m = n$ with the same number of equations and variables
      \item Infinite solutions if $m < n$ since there are more variables than equations
    \end{enumerate}
  \end{definition}

  Because of theorem 3.3, we can determine whether a system has 0, 1, or multiple solutions by looking at the extended matrix's Echelon form. The case for 0 solutions is easy. 

  \begin{theorem}
    The system $A x = b$ has 0 solutions if and only if ref$(\Tilde{A})$ contains a row in the form 
    \begin{equation}
      \begin{pmatrix}
      0 & 0 & \cdots & 0 & c
      \end{pmatrix}, \; c \neq 0
    \end{equation}
  \end{theorem}
  \begin{proof}
    The existence of this row is equivalent to the linear equation
    \begin{equation}
      0 x_1 + 0 x_2 + \cdots + 0 x_n = c, \; c \neq 0
    \end{equation}
    which is absurd and cannot have any solution. Under the hyperplane interpretation, we can visualize all the hyperplanes failing to have a common point. 
  \end{proof}

  \begin{corollary}
    Given $m \times n$ matrix $A$, if $m > n$ and the row vectors of $A$ are all linearly independent, then the system $A x = b$ has 0 solutions. 
  \end{corollary}

  \begin{theorem}
    The system $A x = b$ has 1 solution if and only if ref$(A)$ is diagonal. 
  \end{theorem}

  \begin{proof}
    ref$(A)$ being diagonal implies that there exists at least one solution and also implies the absence of any free variables. 
  \end{proof}

  \begin{theorem}
    The system $A x = b$ has multiple solutions if and only if ref$(A)$ has free variables. 
  \end{theorem}
  \begin{proof}
    Clear. 
  \end{proof}

  \begin{theorem}
    A vector is a solution to the system of equation $A x = b$ if and only if it is of the form 
    \begin{equation}
    a + \text{Null}\; (A)
    \end{equation}
    where $a$ is one solution. 
  \end{theorem}

  \begin{example}
    Given a system of linear equations 
    \begin{align*}
      x + 3 y - 2z = 5 \\
      3 x + 5 y + 6 z = 7 \\
      2 x + 4 y + 3 z = 8
    \end{align*}
    We put it into extended matrix form $A$ and perform Gauss Elimination to get rref$(A)$. 
    \begin{equation}
      \begin{pmatrix}
      1 & 3&-2&5 \\ 3&5&6&7\\ 2&4&3&8
      \end{pmatrix} \rightarrow \begin{pmatrix}
      1&3&-2&5\\ 0&1&-3&2 \\ 0&0&1&2 \end{pmatrix} \rightarrow \begin{pmatrix}
      1&0&0&-15 \\ 0&1&0&8 \\ 0&0&1&2
      \end{pmatrix}
    \end{equation}
    So, rref$(A)$ has the solution $(-15, 8, 2)$ and it is unique because there are no free variables. 
  \end{example}

  This leads to the following theorem. 

  \begin{theorem}
    The set of $n$ linear equations with $n$ variables can be expressed in the form of $A x = b$, where $A$ is an $n \times n$ matrix. 
    \begin{equation}
      A x = b \text{ has a unique solution} \iff A \text{ is nonsingular} \iff \text{rk}\,(A) = n
    \end{equation}
  \end{theorem}
  \begin{proof}
    $A$ is nonsingular is equivalent to saying that rref$(A) = I_n$, where $I_n$ is the $n \times n$ identity matrix. This clearly means that rref$(\Tilde{A})$ will always reveal unique solutions. 
  \end{proof}

\subsection{Numerical Methods}

  In this section, we will concern ourselves with a system of equations with only \textit{one solution}, represented by the matrix equation
  \begin{equation}
    A x = b
  \end{equation}
  where $A$ is an invertible square matrix, $b$ some given vector, and $x$ the vector of unknowns to be determined. 

  An algorithm for solving the system takes as inputs the matrix $A$ and vector $b$ and outputs some approximation to the solution $x$. However, with billions of arithmetic operations on top of each other, the errors can accumulate. Algorithms for which this does not happen are said to be \textbf{arithmetically stable}. 

  The use of finite digit arithmetic places an absolute limitation on the accuracy with which the solution can be determined. To demonstrate this, let us imagine a change $\delta b$ being made in the vector $b$, which causes a corresponding change in $x$, denoted $\delta x$. 
  \begin{equation}
    A(x + \delta x) = b + \delta b \implies A \delta x = \delta b
  \end{equation}
  To compare the changes in $x$ with the changes in $b$, we define the following variable. 

  \begin{definition}
    The \textbf{relative change in x with the relative change in $b$} is the quantity
    \begin{equation}
      \frac{|\delta x|}{|x|} \bigg/ \frac{|\delta b|}{|b|}
    \end{equation}
    where the norm is convenient for the problem (usually a numerical approximation of the Euclidean norm for floating point numbers). We will assume the use of the Euclidean norm from now on. We can rewrite the value as the expression below with the following upper bound, denoted by $\kappa (A)$, called the \textbf{condition number}. 
    \begin{equation}
      \frac{|b|}{|x|} \frac{|\delta x|}{|\delta b|} = \frac{|Ax|}{|x|} \frac{|A^{-1} \delta b|}{|\delta b|} \leq |A||A^{-1}| \equiv \kappa (A)
    \end{equation}
    where $|A|$ is the matrix norm of $A$. 
  \end{definition}

  A high value of this relative change would mean that small perturbations in $b$ would cause large changes in $x$.

  Note that $\kappa(A) \geq 1$. Notice also that the higher the condition number $\kappa (A)$, the harder it is to solve the equation $A x = b$, and $\kappa(A) = \infty$ when $A$ is not invertible. For a $k$-digit floating point arithmetic, the relative error in $b$ can be as large as $10^{-k}$, meaning that the relative error in $x$ can be as large as $10^{-k} \kappa (A)$. 

  Let $\beta$ be the largest absolute value of the eigenvalues of $A$ and $\alpha$ as the smallest absolute value of the eigenvalues of $A$. Then 
  \begin{equation}
    \beta \leq |A|, \frac{1}{\alpha} \leq |A^{-1}| \implies \frac{|\beta|}{|\alpha|} \leq \kappa(A)
  \end{equation}

  An algorithm that generates an exact solution after a finite number of arithmetic steps is called a \textit{direct method} (e.g. Gauss Elimination). An algorithm that generates successive approximations that converge onto the solution is called an \textit{iterative method}. 

  The methods mentioned in this section will be iterative. 

  \begin{definition}
    Let $\{ x_n\}$ be the sequence of approxmations generated by such an algorithm. The deviation of $x_n$ from the true value $x$ is caelled the \textbf{error at the $n$th stage}, denoted by $e_n$. 
    \[e_n \equiv x_n - x\]
    The amount by which the $n$th approximation fails to satisfy the equation $Ax = b$ is called the $n$th residual, denoted by $r_n$. 
    \[r_n \equiv A x_n - b\]
    Error and residual are related by the equation. 
    \[r_n = A e_n\]
  \end{definition}

  Note that since we do not know $x$, the error cannot be calculated, but the residuals can be. We further restrict our scope to solving linear systems in which $A$ is real, positive, and self-adjoint. Clearly, we already know that $|A| = \beta$, and since $A$ is positive, we can conclude that 
  \begin{equation}
    |A^{-1}| = \frac{1}{\alpha}
  \end{equation}

  which implies that
  \begin{equation}
    \kappa(A) = \frac{\beta}{\alpha}
  \end{equation}

\subsubsection{Method of Steepest Descent}

  \begin{theorem}
    Assume that $n \times n$ matrix $A$ is self-adjoint. The solution of $Ax = b$ minimizes the functional 
    \begin{equation}
      E(y) \equiv \frac{1}{2} (y, A y) - (y, b)
    \end{equation}
    where $(\cdot, \cdot)$ is the Euclidean dot product. That is, the solution $x$ is
    \begin{equation}
      x = \min \big\{ E(y) \big\} = \min \Big\{ \frac{1}{2} (y, Ay) - (y, b) \Big\}
    \end{equation}
  \end{theorem}
  \begin{proof}
    Add to $E(y)$ a constant, that is a term independent of $y$ to define a new function $F$. 
    \[F(y) \equiv E(y) + \frac{1}{2} (x, b)\]
    Then, by self adjointness of $A$, we can express $F(y)$ as 
    \begin{align}
      F(y) & = \frac{1}{2} (y, Ay) - (y, b) + \frac{1}{2} (x, b) \\
           & = \frac{1}{2} (y, Ay) - \frac{1}{2} (y, Ax) + \frac{1}{2} (Ax, x) - \frac{1}{2} (Ay, x) \\
           & = \frac{1}{2} \big( y, A(y-x)\big) + \frac{1}{2} \big( A(x-y), x\big) \\
           & = \frac{1}{2} \big( y - x, A(y - x)\big) 
    \end{align}
    Since $F(x) = 0$ and $F(x) \geq 0$ (since it is an inner product with repsect to $y-x$), $F(y)$, and also $E(y)$, takes a minimum at $y =x$. 
  \end{proof}

  Now to actually compute the value of $x$, we us the method of steepest descent. Note that $E: \mathbb{R}^m \longrightarrow \mathbb{R}$, so we can utilize ordinary calculus on it. The gradient of $E$ can be computed by the formula 
  \begin{equation}
    \text{grad}\,E(y) = A y - b
  \end{equation}

  So, if our $n$th approximation is $x_n$, then the $(n+1)$st approximation, $x_{n+1}$, is calculated as
  \begin{equation}
    x_{n+1} = x_n - s (Ax_n - b)
  \end{equation}

  where $s$ is the step length in the direction $-$grad$\,E$. By calculating the residual $r_n = A x_n - b$, we can rewrite the above to
  \begin{equation}
    x_{n+1} = x_n - s r_n
  \end{equation}

  Rather than keeping $s$ constant, we can actually determine an optimal value of $s$ at the $n$th step, denoted $s_n$, which minimizes $E(x_{n+1})$. This quadratic minimum problem is easily solved, since
  \begin{align}
    E(x_{n+1}) & = \frac{1}{2} \big(x_n - s r_n, A(x_n - s r_n) \big) - \big( x_n - s r_n, b\big) \\
    & = E(x_n) - s (r_n, r_n) + \frac{1}{2} s^2 (r_n A r_n)
  \end{align}

  By taking the derivative and computing the value of $s$ where $E (x_{n+1}) = 0$, we see that the minimum is reached when 
  \begin{equation}
    s = s_n \equiv \frac{(r_n, r_n)}{(r_n, A r_n)}
  \end{equation}

  \begin{theorem}
    The sequence of approximations $\{x_n\}$, with $s$ optimized to be $s_n$, converges to the solution of $A x = b$. 
  \end{theorem}

  The error bound for this algorithm is 
  \begin{equation}
    \|e_n\|^2 \leq \frac{2}{\alpha} \bigg( 1-\frac{1}{\kappa(A)} \bigg)^n F(x_0)
  \end{equation}
  which shows that the error $e_n$ tends to $0$ in $\mathbb{R}^m$. However, this algorithm has a very slow rate of convergence for large $\kappa(A)$. 

\subsubsection{Method of Chebyshev Polynomials}

  The disadvantage of the method of steepest descent mentioned in the end of the last subsection renders it quite outdated and obsolete. this next method has a much better error bound that can handle large values of $\kappa$ more efficiently. However, we will need a positive lower bound $m$ for the smallest eigenvalue of $a$ and an upper bound $m$ for the largest eigenvalue. That is, 
  \begin{equation}
    m \leq \alpha, \beta \leq m
  \end{equation}

  and all the eigenvalues of $a$ lie in the interval $[m, m]$. It follows that
  \begin{equation}
    \kappa = \frac{\beta}{\alpha} < \frac{m}{m}
  \end{equation}

  We generate the same sequence of approximations $\{x_n\}$ by the same recursion formula
  \begin{equation}
    x_{n+1} = x_n - s(a x_n - b) \iff x_{n+1} = (i - s_n a) x_n + s_n b
  \end{equation}

  Since the solution of $x$ satisfies the formula; that is, since $x = (i - s_n a) x + s_n b$, we subtract this equation from the top to get
  \begin{equation}
    e_{n+1} = (i - s_n a) e_n
  \end{equation}

  Doing this recursively, we can deduce an explicit formula 
  \begin{equation}
    e_n = p_n (a) e_0 = \prod_{n=1}^n (1 - s_n a)
  \end{equation}

  This allows us to estimate the size of $e_n$. 
  \begin{equation}
    ||e_n|| \leq ||p_n (a)|| ||e_0||
  \end{equation}

  The norm of a self adjoint matrix $a$ is the largest $|a|$, where $a$ is the eigenvalue, and the spectral mapping theorem states that the eigenvalues $p$ of $p_n (a)$ are of the form $p = p_n (a)$, where $a$ is an eigenvalue of $a$. this means that
  \begin{equation}
    ||a|| \leq \max_{m \leq a \leq m} |a| \implies ||p_n (a)|| \leq \max_{m \leq a \leq m} |p_n (a)|
  \end{equation}

  So, we are left with the bound
  \begin{equation}
    ||e_n|| \leq ||e_0|| \max_{m \leq a \leq m} |p_n (a)|
  \end{equation}

  To get the best estimate of $e_n$, we have to choose the $s_1, s_2, \cdots, s_n$ so that the polynomial $p_n$ has a small maximum on the interval $[m, m]$. note that the polynomial $p_n$ satisfies the normalizing condition 
  \begin{equation}
    p_n(0) = 1
  \end{equation}

  Oo find such a polynomial, we must first define chebyshev polynomials. 

  \begin{definition}
    The \textbf{$n$th chebyshev polynomial} $t_n$ is defined for $-1 \leq u \leq 1$ by
    \begin{equation}
      t_n (u) = \cos (n \theta), \;\; u = \cos(\theta)
    \end{equation}
  \end{definition}

  \begin{theorem}
    Among all polynomials $p_n$ of degree $n$ that satisfy $p_n (0) = 1$, the one that has the smallest maximum on $[m, m]$ is the \textit{rescaled chebyshev polynomial} that rescales values from $[-1, 1]$ to the interval $[m, m]$ while preserving the condition that $p_n (0) = 1$. this polynomial is expressed as
    \begin{equation}
      p_n (a) \equiv t_n \bigg( \frac{m+m-2a}{m-m} \bigg) \bigg/ t_n \bigg(\frac{m+m}{m-m} \bigg)
    \end{equation}
  \end{theorem} 

  Now, assuming that $m/m \approx \kappa$, 
  \begin{equation}
    t_n \bigg( \frac{m+m}{m-m} \bigg) = t_n \bigg(\frac{\frac{m}{m} + 1}{\frac{m}{m}-1} \bigg) \approx t_n \bigg(\frac{\kappa + 1}{\kappa - 1} \bigg)
  \end{equation}

  since $|t_n (u)| \leq 1$ for $|u| \leq 1$, this also implies that
  \begin{equation}
    t_n \bigg( \frac{m + m - 2a}{m-m} \bigg) \leq 1
  \end{equation}

  Combining this together, we get
  \begin{equation}
    ||e_n|| \leq ||e_0|| \max_{m \leq a \leq m} |p_n (a)| = ||e_0|| \bigg/ t_n \bigg( \frac{\kappa+1}{\kappa-1} \bigg)
  \end{equation}

  It is a fact that higher order chebyshev polynomials tend to infinity faster once the value reaches out of $[-1, 1]$, meaning that as $n \rightarrow \infty$, $t_n\big( (\kappa+1)/(\kappa-1)\big)$ will also tend to infinity (note that $(\kappa+1)/(\kappa-1)$ is a constant, implying that $e_n$ tends to $0$ as $n$ tends to infinity. the error bound for $e_n$ is given by the following 
  \begin{equation}
    ||e_n|| \leq 2 \bigg( 1 + \frac{2}{\sqrt{\kappa}} \bigg)^{-n} ||e_0|| \approx 2 \bigg( 1 - \frac{2}{\sqrt{\kappa}} \bigg)^{n} ||e_0||
  \end{equation}

  Once again, this confirms that $e_n \rightarrow 0$ as $n \rightarrow \infty$. furthermore, when $\kappa$ is large, the error bound works with $\sqrt{\kappa}$, which is must smaller than $\kappa$ itself. so, $e_n$ converges much faster through this algorithm than through the method of steepest descent. 


